% §3 System Architecture — draft (Paper I, V19 phase 1)
\section{System Architecture}
\label{sec:system}

We target a cloud--edge--device stack for robotic tongue-image acquisition in traditional Chinese medicine (TCM) tele-diagnosis.
The design separates \emph{real-time} acquisition and quality control at the edge from \emph{offline} or high-latency cloud analytics, while keeping algorithm modules testable outside the robot operating system (ROS) runtime.

\subsection{Device Layer}
The device layer consists of a Franka Research 3 (FR3) manipulator, RGB imaging, and ROS~2-based motion execution.
MoveIt provides point-to-point (PTP) joint-space motions.
For Paper~I we expose two named \emph{skills}---\texttt{go\_to\_tongue\_pose} and \texttt{go\_to\_face\_pose}---stored as calibrated joint configurations in \texttt{poses.yaml} and invoked through the edge gateway.
This avoids impedance or stiffness tuning interfaces in the paper experiments (\texttt{PAPER1\_MODE=1} disables those routes).
Simulated and fake-hardware stacks (MoveIt fake controller, optional Gazebo camera) support reproducible auxiliary trials without a physical robot.

\subsection{Edge Layer}
The edge layer is the primary contribution surface for closed-loop acquisition.
A FastAPI server (\texttt{franka\_api\_server}) acts as the edge gateway on port~8000, bridging HTTP clients to ROS~2 actions for motion and error recovery.
\textbf{Edge-IQA} scores each captured frame with a lightweight quality index $Q_{\mathrm{img}}$ (sharpness and exposure fusion), executed in the \texttt{doctor/paper1} tree within the \texttt{franka\_ros2} monorepo and mounted at the gateway via subprocess in phase~2.
\textbf{LangGraph} implements confidence-aware routing over a small state machine: \texttt{capture} $\rightarrow$ \texttt{edge\_iqa} $\rightarrow$ \texttt{route}, with branches \texttt{upload\_cloud}, \texttt{resample\_edge}, and \texttt{fail\_safe}.
The HTTP client \texttt{sim/franka\_bridge} connects LangGraph nodes to gateway endpoints (\texttt{/vision/evaluate}, \texttt{/motion/skills/...}).
Phase~1 uses a placeholder vision endpoint returning fixed $Q_{\mathrm{img}}$ for API contract validation.

\subsection{Cloud Layer}
The cloud layer hosts higher-level vision models and EMR integration (stub in Paper~I).
Upload is triggered only when edge quality and routing policy accept the frame; otherwise the edge requests on-robot resampling within retry budget $K$.
This division keeps the 1\,kHz ROS control loop free of learning-based perception.
Edge-IQA never writes directly to the hardware loop; it selects among pre-calibrated skills or a fail-safe return path exposed by the gateway.
Thus a perception fault degrades the acquisition workflow, not the low-level controller.

\subsection{Data and Timing Artifacts}
Each closed-loop trial logs ISO~8601 timestamps (\texttt{t\_capture}, \texttt{t\_iqa}, \texttt{t\_route}) and decisions to JSONL for supplementary reproducibility.
Main quantitative results (RTT percentiles, valid-frame rate, retry rate) are produced \emph{offline} by \texttt{run\_matrix\_tcm.py} on ShezhenV3-COCO test images (baselines B0--B3); the franka auxiliary track (M6) reports trend-consistent RTT only, not primary table numbers.
The same schema is used for offline replay and the auxiliary Franka track, so every uploaded, retried, or aborted frame can be traced to a score, flag set, retry count, and timestamp.

\subsection{Threats to Validity}
Offline RTT aggregates a documented latency model rather than live WAN traces; relative baseline ordering is the primary claim.
ShezhenV3 blur labels are synthetically injected on validation/test frames to emulate motion defocus when human quality annotations are unavailable.
Edge-IQA overlap between clear and blurred real cohorts (M5 $\rho{=}0.47$) motivates future ROI cropping; the routing policy remains interpretable and deployable on CPU gateways.

\subsection{Implementation Mapping}
Table~\ref{tab:impl-map} summarizes the mapping from architecture layers to repositories.

\begin{table}[t]
\centering
\caption{Layer-to-repository mapping (Paper~I).}
\label{tab:impl-map}
\small
\resizebox{\columnwidth}{!}{%
\begin{tabular}{lll}
\hline
Layer & Component & Location \\
\hline
Device & MoveIt, skills & \texttt{franka\_ros2} \\
Edge & FastAPI gateway & \texttt{franka\_api\_server} \\
Edge & Edge-IQA, LangGraph & \texttt{doctor/paper1} \\
Cloud & Vision stub & \texttt{doctor/paper1} (stub) \\
\hline
\end{tabular}%
}
\end{table}

\begin{figure*}[t]
  \centering
  \includegraphics[width=\linewidth]{fig1_system_overview.pdf}
  \caption{Cloud--edge--device stack for robotic tongue-image acquisition (Paper~I).}
  \label{fig:system-overview}
\end{figure*}

Figure~\ref{fig:system-overview} shows the end-to-end information flow; device motions are commanded only through the edge gateway, and IQA code is not embedded inside ROS hardware plugins.
